\chapter{Giới thiệu}
\label{ch:intro}

\section{Bối cảnh}

Trong gần hai thập kỷ trở lại đây, tốc độ xung nhịp của bộ xử lý gần như dừng lại ở mức
3--5~GHz do giới hạn về công suất tiêu thụ và tản nhiệt (hiện tượng \emph{power wall}).
Thay vì tăng tốc một nhân, các nhà sản xuất chuyển sang tích hợp nhiều nhân trên cùng một chip:
ngay cả một máy tính xách tay phổ thông ngày nay cũng có từ 8 đến 16 luồng phần cứng, kèm theo
các đơn vị tính toán vector (SIMD) rộng 256 hoặc 512 bit~\cite{hennessy2019}. Hệ quả là hiệu
năng phần mềm không còn ``tự động'' tăng theo thế hệ phần cứng; muốn tận dụng tài nguyên tính
toán, lập trình viên buộc phải \emph{song song hóa} chương trình.

Song song hóa dễ nhất với các bài toán mà mỗi phần tử dữ liệu được xử lý độc lập (cộng hai
vector, áp bộ lọc lên từng điểm ảnh) -- loại bài toán thường gọi là \emph{embarrassingly
parallel}. Tuy nhiên, rất nhiều thuật toán quan trọng chứa \emph{phụ thuộc dữ liệu} giữa các
bước lặp, khiến việc chia việc cho nhiều luồng không còn hiển nhiên. Ví dụ điển hình, đồng
thời là chủ đề của báo cáo này, là phép \textbf{tính tổng tiền tố} (\emph{prefix sum},
hay \emph{scan}):
\begin{equation}
  y_i \;=\; x_0 \oplus x_1 \oplus \cdots \oplus x_i, \qquad i = 0, 1, \ldots, n-1,
  \label{eq:intro-scan}
\end{equation}
trong đó \(\oplus\) là phép toán hai ngôi có tính kết hợp (thường là phép cộng). Cài đặt tuần
tự của~\eqref{eq:intro-scan} chỉ là vòng lặp \(y_i = y_{i-1} \oplus x_i\); mỗi bước phụ thuộc
trực tiếp vào bước liền trước. Thoạt nhìn, phép toán này có vẻ \emph{bản chất tuần tự}.

Blelloch~\cite{blelloch1990} đã chỉ ra rằng điều đó không đúng: tổng tiền tố là một ví dụ
điển hình về phép tính trông có vẻ tuần tự nhưng lại có thuật toán song song hiệu quả. Hơn
nữa, scan được xem là một trong những \emph{khối xây dựng} (building block) cơ bản nhất của
thuật toán song song, có vai trò tương tự phép duyệt mảng trong thuật toán tuần tự. Rất nhiều
bài toán -- sắp xếp cơ số, sắp xếp nhanh, lọc phần tử (stream compaction), cấp phát bộ nhớ
động, phân tích từ vựng, giải hệ truy hồi, dựng histogram, bảng tổng vùng trong xử lý ảnh --
có thể quy về một hoặc vài lần gọi scan~\cite{blelloch1990,harris2007,sengupta2007}. Vì vậy,
một cài đặt scan song song nhanh có ý nghĩa thực tiễn lớn: nó trực tiếp tăng tốc cả một lớp
ứng dụng.

Phần lớn tài liệu kinh điển về scan song song tập trung vào mô hình lý thuyết
PRAM~\cite{blelloch1990,ladner1980} hoặc vào GPU với hàng nghìn
luồng~\cite{harris2007,merrill2016}. Trên CPU đa nhân -- nơi số luồng chỉ khoảng chục, có hệ
thống bộ nhớ đệm nhiều cấp và băng thông bộ nhớ hạn chế -- các đánh đổi lại rất khác: một thuật
toán ``tối ưu'' trên lý thuyết chưa chắc đã nhanh trong thực tế. Đây chính là động lực của đề
tài: \emph{cài đặt các thuật toán scan song song bằng OpenMP trên CPU đa nhân và đánh giá hiệu
năng một cách có hệ thống}.

\section{Phát biểu bài toán}

\paragraph{Đầu vào.} Mảng \(x = [x_0, x_1, \ldots, x_{n-1}]\) gồm \(n\) số nguyên 64 bit lưu
liên tiếp trong bộ nhớ, và số luồng \(p\).

\paragraph{Đầu ra.} Mảng \(y\) cùng kích thước thỏa mãn~\eqref{eq:intro-scan} với \(\oplus\)
là phép cộng (\emph{inclusive scan}). Kết quả song song phải \textbf{trùng khớp tuyệt đối} với
kết quả tuần tự.

\paragraph{Yêu cầu.}
\begin{enumerate}
  \item Cài đặt bằng C và OpenMP các thuật toán scan song song tiêu biểu: Hillis--Steele,
        Blelloch (up-sweep/down-sweep), scan chia khối và chỉ thị \code{scan} của OpenMP 5.0.
  \item Kiểm chứng tính đúng trên nhiều kích thước dữ liệu và số luồng.
  \item Đánh giá hiệu năng: thời gian chạy, speedup, hiệu suất, khả năng mở rộng (strong/weak
        scaling), băng thông bộ nhớ đạt được; phân tích nguyên nhân của các kết quả.
\end{enumerate}

Việc chọn số nguyên 64 bit thay vì số thực là có chủ đích: phép cộng dấu phẩy động
\emph{không} có tính kết hợp tuyệt đối, nên thứ tự cộng khác nhau giữa các thuật toán sẽ cho
sai khác nhỏ, gây khó cho kiểm thử. Với số nguyên, mọi sai khác dù chỉ một bit đều là lỗi.

\section{Mục tiêu và phạm vi}

\begin{itemize}
  \item \textbf{Lý thuyết:} trình bày định nghĩa, các thuật toán scan song song kinh điển cùng
        phân tích theo mô hình work--span; làm rõ khái niệm \emph{work-efficiency} và vì sao nó
        quan trọng khi số luồng nhỏ hơn nhiều so với kích thước dữ liệu.
  \item \textbf{Cài đặt:} xây dựng một thư viện nhỏ gồm sáu phiên bản scan với giao diện thống
        nhất, kèm bộ kiểm thử tự động và chương trình đo hiệu năng.
  \item \textbf{Thực nghiệm:} đo đạc trên một CPU lai (hybrid) hiện đại, đối chiếu với giới hạn
        băng thông bộ nhớ để trả lời câu hỏi: \emph{scan song song trên CPU nhanh được đến đâu
        và bị giới hạn bởi điều gì?}
  \item \textbf{Ứng dụng:} minh họa dùng scan để song song hóa bài toán lọc phần tử.
\end{itemize}

\noindent\textbf{Phạm vi.} Đề tài giới hạn trong mô hình bộ nhớ chia sẻ trên một nút tính toán,
sử dụng OpenMP. GPU (CUDA) và bộ nhớ phân tán (MPI) chỉ được đề cập ở mức tổng quan.

\section{Đóng góp của nhóm}

\begin{enumerate}
  \item Cài đặt và kiểm thử sáu phiên bản scan bằng OpenMP, trong đó phiên bản tối ưu
        \emph{Blocked-opt} kết hợp chia tile vừa bộ nhớ đệm L2, chống false sharing và vector
        hóa, giảm lưu lượng bộ nhớ từ 24 xuống 16 byte cho mỗi phần tử.
  \item Xây dựng quy trình đo hiệu năng có thể tái lập (script tự động, nhiều lần lặp, lấy
        trung vị, ghi lại cấu hình máy) và công cụ sinh biểu đồ, bảng LaTeX tự động.
  \item Phân tích kết quả dựa trên mô hình roofline: với dữ liệu lớn, scan là bài toán
        \emph{giới hạn bởi băng thông bộ nhớ}, nên speedup bị chặn bởi tỉ số băng thông đa
        luồng/đơn luồng chứ không phải bởi số nhân.
  \item Ghi nhận các vấn đề thực tế: thuật toán ``đẹp'' trên lý thuyết (Hillis--Steele,
        Blelloch) chậm hơn cả bản tuần tự trên CPU; chỉ thị \code{omp scan} của GCC lỗi khi
        vòng lặp rỗng; ảnh hưởng của kiến trúc lai P-core/E-core tới cân bằng tải.
\end{enumerate}

\section{Cấu trúc báo cáo}

\Cref{ch:background} trình bày cơ sở lý thuyết: định nghĩa scan, mô hình work--span, các định
luật giới hạn tăng tốc, mô hình roofline, chuẩn OpenMP và đặc điểm kiến trúc CPU.
\Cref{ch:algorithms} trình bày chi tiết các thuật toán scan song song, phân tích độ phức tạp và
chứng minh tính đúng. \Cref{ch:impl} mô tả cài đặt bằng OpenMP, kỹ thuật tối ưu và kiểm thử.
\Cref{ch:experiments} trình bày môi trường, phương pháp và kết quả thực nghiệm kèm phân tích.
\Cref{ch:conclusion} tổng kết, nêu hạn chế và hướng phát triển.
